% \foreach 的局部作用域问题是 TeX “组”（Group）机制的直接体现，
% 而非现代编程语言中的“沙盒”概念。每次循环迭代都相当于进入了一个
% 临时的 TeX 组，组内的宏定义和变量修改在组结束（即本次迭代结束）
% 时会被自动还原。针对在循环中计算出某个值（例如动态生成的节点名、
% 坐标或数组元素），并希望在循环结束后继续使用它的需求，使用 \xdef
% （或 \global\edef，或 \global\let）进行全局赋值是最经典且
% 有效的解决方案。

% 方法：将正多边形投影到任意平面上

\documentclass{standalone}
\usepackage{tikzplot}

\begin{document}

\begin{tikzpicture}
  \tikzmath{
    % 正 n 边形的边数
    \size = 5;
    % 外接圆和内切圆的圆心
    \O{1,1} = 0; \O{2,1} = 0; \O{3,1} = 1; 
    \P{1,1} = 1;
    \P{2,1} = 1/4;
    \P{3,1} = 1;
    % 外接圆的半径
    \ra = 2;
    % 对角线形成的正 n 边形内切圆的半径
    \rb = \ra * cos(360/\size);
    % center of projection
    \Center{1} = 0; \Center{2} = 0; \Center{3} = 6;
    % source plane
    \S{1} = 0; \S{2} = 0; \S{3} = 0;
    \m{1} = 0; \m{2} = 0; \m{3} = 1;
    % target plane
    \T{1} = 0; \T{2} = 0; \T{3} = 0;
    \n{1} = -1; \n{2} = -3/8; \n{3} = 1/2; 
  }
  
  \tikzset{
    % 以 O 为圆心 radius 为半径的圆
    conics/define/circle/center-radius={\O,\ra,\Circlea},
    conics/define/circle/center-radius={\O,\rb,\Circleb},
    space/perspective transform={\Center,\S,\m,\T,\n,\H},
    mat/inverse={\H,3,\Hi},
    mat/congruence={\Circlea,\Hi,3,\CCa},
    mat/congruence={\Circleb,\Hi,3,\CCb},
    mat/normalize={\CCa,3,3,\Ca},
    mat/normalize={\CCb,3,3,\Cb},
    conics/intersect cl={\Ca,\O,\P,\Qa,\Qb},
    conics/dehomogenize={\Qa,\QQa},
  }

  \tikzmath{
    % 正 n 边的起始点, 每次都修改, 用于排除已找到顶点
    \PUprev = \Qa{1,1};
    \PVprev = \Qa{2,1};
    \PWprev = \Qa{3,1};
    % 正 n 边与内切圆的切点, 每次都修改, 用于排除已找到的边
    \TUprev = 0;
    \TVprev = 0;
    \TWprev = 1;
  }
  
  \axes[grid] (-5:5,-5:5);
  % \conic[draw=magenta] (\Circlea);
  % \conic[draw=magenta] (\Circleb);
  \conic[thick,draw=teal] (\Ca);
  \conic[thick,draw=navy] (\Cb);

  \foreach \i in {1,...,\size}{%sandbox
    \tikzmath{
      % 当前的顶点和切点
      \Pcurr{1,1}=\PUprev; \Pcurr{2,1}=\PVprev; \Pcurr{3,1}=\PWprev;
      \Tcurr{1,1}=\TUprev; \Tcurr{2,1}=\TVprev; \Tcurr{3,1}=\TWprev;
    }
    
    \tikzset{
      mat/stdout={\Pcurr,3,1},
      mat/stdout={\Tcurr,3,1},
      conics/tangents={\Cb,\Pcurr,\Ta,\Tb},
      conics/dehomogenize={\Pcurr,\PPa},
      conics/exclude={\Ta,\Tb,\Tcurr,\next},
      conics/intersect cl={\Ca,\Pcurr,\next,\P,\Q},
      conics/exclude={\P,\Q,\Pcurr,\Pnext},
      conics/dehomogenize={\Pnext,\PPn},
      conics/dehomogenize={\next,\TTn},
    }
    % 导出下一个顶点和切点
    \xdef\PUprev{\Pnext{1,1}}
    \xdef\PVprev{\Pnext{2,1}}
    \xdef\PWprev{\Pnext{3,1}}
    \xdef\TUprev{\next{1,1}}
    \xdef\TVprev{\next{2,1}}
    \xdef\TWprev{\next{3,1}}
 
    \coordinate (P) at (\PPa{1},\PPa{2});
    \coordinate (Q) at (\PPn{1},\PPn{2});
    \draw[red] (P) -- (Q);
    \foreach \p/\placement in {P/above}{
      \fill[red] (\p) circle (1.5pt);
     \draw (\p) node[\placement] {$\p\i$};
    }
  }
\end{tikzpicture}

\end{document}
